\Chapter{Dane i przygotowanie zbioru}
\label{ch:data}

\section{Źródło, reprezentacja i jakość}

Źródłem danych jest FastF1, w audycie użyte w wersji \code{3.9.0a0}.
Człowiek pobrał dane do lokalnej pamięci podręcznej z użyciem własnego
uwierzytelnienia F1 TV; dalszy eksport i analiza działają offline.
Zespół może więc korzystać z tej samej utrwalonej wersji bez ponownego
pobierania i bez przekazywania poświadczeń. Obok okrążeń zachowano
statusy toru, komunikaty kontroli wyścigu i zapowiedzi liczby okrążeń.
Zestawienia wizyt Jolpica posłużyły do kontroli zgodności zdarzeń.
Pogoda była sprawdzana w audycie, ale nie jest cechą ani składnikiem
bieżącego eksportu; telemetria o wysokiej częstotliwości pozostaje
poza zakresem.

Kandydat danych zawiera 92 żądane i 92 załadowane wyścigi oraz
101\,290 surowych wierszy. Są to liczby pokrycia wejścia, nie rozmiar
gotowego zbioru uczącego. Każdy wiersz opisuje okrążenie jednego
kierowcy: jego czas zakończenia i trwania, pozycję, status toru,
znaczniki wjazdu i wyjazdu oraz historię opon. \code{Compound} oznacza
mieszankę, \code{TyreLife} wiek opony, a \code{Stint} nominalny numer
odcinka historii opon. Numer kierowcy interpretuje się w kontekście
wyścigu, nie jako stały identyfikator osoby między sezonami.

Czasy okrążeń i znaczników pit stopu zapisuje się w sekundach sesji.
Brak znacznika wjazdu nie oznacza czasu równego zero. Nieznane pomiary
pozostają puste; tekstowe napisy \code{nan} i \code{None} w polach
mieszanki i statusu są normalizowane do rzeczywistych braków.
Poprawka ta ujawniła 780 braków mieszanki wcześniej błędnie liczonych
jako obecne wartości. Nie odtwarza ona jednak historii opon.
Sam sukces odczytu sesji nie gwarantuje, że każdą jej obserwację da się
wykorzystać w uczeniu~\cite{findings}.

Przyjęta polityka wyklucza cały wyścig, jeśli występuje choć jedno
okrążenie na oponach INTERMEDIATE/WET lub status zawierający kod
czerwonej flagi. Audyt wskazał 18 takich wyścigów. Pozostają okrążenia
z samochodem bezpieczeństwa (SC) i wirtualnym samochodem bezpieczeństwa
(VSC). Dla pary wyścig--kierowca powodem wykluczenia jest przerwa większa
niż trzy brakujące okrążenia lub ponad 10\% braków mieszanki bądź wieku
opony. Szczególne reguły opon w Katarze 2023, Monako 2025 i Katarze 2025
są osobno oznaczane. Eksporter zapisuje powody wykluczeń jako metadane,
nie usuwa materiału surowego. Docelowe zastosowanie reguł przy budowie
próbek musi zachować tę możliwość kontroli.

\section{Znaczenie etykiety}

Pierwotna reguła przyjmowała odpowiedź dodatnią, gdy na okrążeniu $t+1$
zapisano \code{PitInTime}, istniało $t+2$, a numer stintu między tymi
okrążeniami wzrósł. Audyt wykazał taki wzrost także przy obsłudze kar.
W przykładzie z tabeli~\ref{tab:penalty} komunikaty kontroli wyścigu
wskazują karę stop/go, a wiek opony rośnie bez przerwania historii.
Oznacza to, że wzrost \code{Stint} nie stanowi samodzielnego dowodu
wymiany opon~\cite{findings}.

\begin{table}[htbp]
\centering\small
\caption{Przejazd związany z karą: Baku 2022, kierowca 6.}
\label{tab:penalty}
\begin{tabular}{rrrll}
\toprule
Okrążenie & Stint & TyreLife & Znacznik alei & Compound \\
\midrule
3 & 1 & 3 & wjazd & MEDIUM \\
4 & 2 & 4 & wyjazd & MEDIUM \\
\bottomrule
\end{tabular}
\end{table}

Niezmieniona mieszanka także nie dowodzi braku wymiany: można założyć
inny komplet tej samej mieszanki. Wiek opony wymaga uwzględnienia
używanych kompletów i braków danych. Nie zatwierdzono zastępczego
klasyfikatora opartego wyłącznie na tych polach. Reguła etykiet pozostaje
zatem otwarta przed ich implementacją. Zmiana odpowiedzi na samą wizytę
w alei oznaczałaby zmianę celu projektu.

Porównanie pomocniczego wskaźnika wjazdu i wzrostu stintu z Jolpica,
z tolerancją jednego okrążenia, dało 3339 dopasowań do 3363 wizyt.
Zgodność ta dotyczy rejestrów przejazdów, nie potwierdzonych wymian opon:
kary mogą występować w obu źródłach. Pozostałych wizyt nie wolno
bez wyjaśnienia zamienić na etykiety ujemne. Podobnie oznaczenia
\code{PitInTime} i \code{PitOutTime} zachowują przypisanie do
źródłowych okrążeń, nawet gdy występuje niewielka rozbieżność ich czasów
względem granic okrążenia.

\section{Dostępność informacji w czasie}

Wyciekiem informacji byłoby użycie przy prognozie obserwacji, która
nie istniała jeszcze w chwili $T$. Kierowcy kończą okrążenia w różnych
momentach, dlatego dla danych rywala sam numer okrążenia nie wystarcza.
Samochód z tyłu może zakończyć to samo okrążenie dopiero po chwili
prognozy. Projekt przewiduje używanie wcześniejszego okrążenia przy
luce do samochodu z tyłu i odcięcie czasowe przy zliczaniu wjazdów
rywali. Każdy strumień i każdy kierowca muszą podlegać temu samemu
ograniczeniu $\mathrm{czas}\leq T$.

Planowana liczba okrążeń również może się zmieniać. Loader FastF1
zachowuje ostatnią niepustą zapowiedź \code{TotalLaps}. W dwóch
zachowanych po wykluczeniach przykładach, \code{2023_14_R}
i \code{2025_11_R}, ostatnie zapowiedzi poprzedzały jednak wszystkie
zarejestrowane końce okrążeń. Wykazano więc zmienność metadanych,
a nie wyciek przyszłej liczby do konkretnych próbek tych dwóch
wyścigów. Zachowanie zapowiedzi z czasem umożliwia dalszą kontrolę,
lecz kontrakt ich użycia i postępowanie przed pierwszą znaną zapowiedzią
pozostają nieustalone. Liczby planowanej nie wolno zastępować
największym numerem okrążenia zaobserwowanym po zakończeniu wyścigu.

Flaga \code{IsAccurate} opisuje przydatność pomiaru czasu, nie
przydatność próbki do etykietowania. Usunięcie wszystkich wierszy
z flagą fałszywą usunęłoby także okrążenia wjazdowe i wyjazdowe.
Flaga ma więc maskować tylko cechy czasowe. Audyt potwierdził również,
że niedokładność pierwszego okrążenia po statusie równym dokładnie
\code{4} nie uogólnia się na każdy status mieszany zawierający
\code{4}~\cite{findings}. Źródłowy status musi zachować to rozróżnienie.

Docelowe cechy i okna wymagają dopiero testów przyczynowych:
ponowne obliczenie po usunięciu wszystkich danych późniejszych niż
$T$ powinno dawać ten sam wynik co obliczenie na pełnym wejściu.
Wypełnianie braków może korzystać wyłącznie z wcześniejszych obserwacji,
a statystyki standaryzacji tylko z części uczącej. Istniejące kontrole
formatu i zachowanie znaczników czasu stanowią do tego podstawę,
ale nie zastępują takiego testu.
